\documentclass{situationist-preprint}

\title{The Branch-Line Case}
\msid{PID-02-b}
\msfamily{PID-02}
\msdiscipline{Personal Identity and the Self}
\mstype{Companion manuscript}
\mskeywords{teletransportation; Branch-Line; branching; what matters; Parfit; extrinsic survival}

\begin{document}
\maketitle

\begin{abstract}
In Parfit's Branch-Line case a scanner records a person without destroying her, a replica is built on Mars, and the original, on Earth, is to die within days. PID-02-a holds that she does not survive as the replica, and its revised statement extends the verdict to every case in which the original exists at some time after the replica is activated. We examine the case in detail. We argue that the verdict on identity is correct, that the replica nevertheless has nearly all of what matters to her, and that the apparent arbitrariness of a verdict that turns on minutes or seconds is the arbitrariness of identity, not of what matters. The dying original has little reason to regard her approaching death as she would an ordinary one.
\end{abstract}

\section{The case and the verdict}

In Branch-Line \citep[ch.~10]{Parfit1984}, the new scanner records the exact state of a person's body and brain without destroying them. A replica is built on Mars. The person on Earth learns that the scan has damaged her heart and that she will die within a few days. She may even speak with her replica by videophone.

PID-02-a holds that a person survives teletransportation if the original is destroyed before the replica is activated, and not otherwise (Theorem~3.3). In Branch-Line the original exists for days after activation, so the process branches and the person does not survive as the replica (Proposition~3.2). She continues on Earth as the original, by the bodily tie-breaker of PID-01, until she dies.

The revision of PID-02-a makes explicit a family of intermediate cases. Call a \emph{delayed destruction} case one in which the original is destroyed $d$ seconds after the replica is activated. For every $d>0$ the process branches, and the person does not survive as the replica. Only where destruction precedes activation does she survive.

\section{What she loses}

\begin{proposition}\label{p:loss}
In Branch-Line, the replica stands to the original in a relation that has nearly all of what matters in survival, though it is not identity.
\end{proposition}

\begin{proof}
At activation the replica is psychologically continuous with the original and as strongly connected with her as she is with herself on the following morning. Whether a later being has what matters depends on the intrinsic character of its relation to the original, not on whether another being stands in a similar relation (PID-03-a, Assumption~2.6). The relation of the replica to the original is intrinsically the same as in Simple Teletransportation, in which it is survival. So the replica has what matters to the degree that a survivor would, except for whatever weight attaches to identity itself; and by PID-06-a, Corollary~3.4, that weight is at most a small $\varepsilon$.
\end{proof}

The original's death on Earth is therefore not like an ordinary death. In an ordinary death nothing in the future stands to the person in a relation that carries what matters. In Branch-Line, something does: a person on Mars who will carry on her projects, love the people she loves and remember her life up to the scan. What she loses is identity, which carries little, and the experiences of her last days, which the replica will not remember. The latter loss is a real one, of the kind suffered by someone who knows she will lose the memory of a short period; it is not the loss of a future.

\section{Seconds and survival}

The delayed destruction cases invite the charge of arbitrariness. If destruction one second before activation is survival and destruction one second after is not, survival seems to depend on a fact too slight to bear the weight.

\begin{objection}
Nothing of importance happens in a second. A theory on which a second makes the difference between survival and death must be wrong.
\end{objection}

\begin{reply}
The charge confuses identity with what matters. PID-02-a, Corollary~3.4, accepts that identity is extrinsic and depends on timing. What matters does not: by Proposition~\ref{p:loss}, the replica's relation to the original has nearly all of what matters for every value of $d$, positive or negative. A theory on which a second made the difference between having and lacking what matters would be absurd; the present theory says that it makes the difference only to a relation that carries little weight. The verdict is sharp because identity is all-or-nothing and one-one, and a one-one relation grounded in a relation that can hold one-many must draw a sharp line where the many begin.
\end{reply}

\begin{objection}
In Branch-Line the original rationally fears her death on Earth. If the replica had nearly everything that matters, the fear would be irrational.
\end{objection}

\begin{reply}
Much of her fear is rational. She has reason to dread the pain of her last days, which the replica will not share, and the experience of dying; fear of one's own future pain does not wait on a theory of identity \citep{Williams1970}. What she has less reason to fear is that death will end everything she cares about in her own future. The residual concern, directed at being the one who continues, corresponds to the small weight of identity. The intensity of the ordinary fear of death is calibrated to cases in which identity and what matters are lost together, and \citet[ch.~12]{Parfit1984} argues that it is miscalibrated in this case; we agree, with the qualification that identity keeps a small residual weight.
\end{reply}

\section{Conclusion}

Branch-Line confirms the extrinsic verdict of PID-02-a. The person on Earth is the original; the replica is a new person; and the line between survival and non-survival falls at the moment of activation. Because what matters does not follow identity across that line, the verdict is not the cruelty it may seem: the person who dies on Earth leaves behind someone who has nearly all of what she had reason to care about in survival.

\begin{thebibliography}{99}\small
\bibitem[Parfit(1984)]{Parfit1984} Parfit, D. (1984). \emph{Reasons and Persons}. Oxford University Press.
\bibitem[Williams(1970)]{Williams1970} Williams, B. (1970). The self and the future. \emph{The Philosophical Review}, 79(2), 161--180.
\end{thebibliography}

\end{document}
